% Unit 039 English translation of chapter3.tex lines 1293--1586.
% Source: Methods of Algebra, Volume 2, commit
% 9a5803ff77dd3257484cb177f851a73770a59dd3 (CC BY 4.0).
% stable-unit-id: o014.aljabr2.chapter3.exercises-hochschild-homology-and-cohomology
% Coverage: Section 3.8 in full.

% segment-id: o014.li2.chapter3.exercises-hochschild-homology-and-cohomology.h001
\section{Exercises: Hochschild Homology and Cohomology}\label{sec:HH}
\hypertarget{o014.aljabr2.chapter3.exercises-hochschild-homology-and-cohomology}{ }

% segment-id: o014.li2.chapter3.exercises-hochschild-homology-and-cohomology.p001
This section returns to concrete mathematics. Let $\Bbbk$ be a commutative
ring. Tensor products $\dotimes{\Bbbk}$ of $\Bbbk$-modules will be abbreviated
to $\otimes$, the $n$-th tensor power of a $\Bbbk$-module $M$ will be written
$M^{\otimes n}$, and we set $M^{\otimes 0}:=\Bbbk$. For a $\Bbbk$-algebra
$R$, as usual, the left and right actions of $\Bbbk$ on every $(R,R)$-bimodule
are assumed to agree. Thus an $(R,R)$-bimodule is also identified with a left
$R\otimes R^{\opp}$-module.

% segment-id: o014.li2.chapter3.exercises-hochschild-homology-and-cohomology.de001
\begin{definition}[Bar complex]\label{def:bar-Hochschild}
	\index{bar complex}
	\index[sym1]{BR@$\mathsf{B}R$}
	Let $R$ be a $\Bbbk$-algebra. For every $n\in\Z_{\geq 0}$, define the
	$(R,R)$-bimodule
% segment-id: o014.li2.chapter3.exercises-hochschild-homology-and-cohomology.q001
	\begin{gather*}
		\mathsf{B}_n R := R \otimes R^{\otimes n} \otimes R = R^{\otimes (n+2)}, \\
		r(r_0 \otimes \cdots \otimes r_{n+1})r'
		= rr_0 \otimes \cdots \otimes r_{n+1} r',
	\end{gather*}
	where $r,r',r_0,\ldots,r_{n+1}\in R$. For every $n\geq1$, define the
	bimodule homomorphism
% segment-id: o014.li2.chapter3.exercises-hochschild-homology-and-cohomology.q002
	\begin{equation*}\begin{split}
		b_n: \mathsf{B}_n R & \to \mathsf{B}_{n-1} R, \\
		b_n \left( r_0 \otimes \cdots \otimes r_{n+1} \right)
		& := \sum_{k=0}^n (-1)^k
		\cdots \otimes r_k r_{k+1} \otimes \cdots .
	\end{split}\end{equation*}
	We call $\mathsf{B}R:=\left(\mathsf{B}_nR,b_n\right)_{n\geq0}$ the
	\emph{bar complex} of $R$; it is a chain complex in the sense of
	Remark~\ref{rem:cochain-vs-chain}.
\end{definition}

% segment-id: o014.li2.chapter3.exercises-hochschild-homology-and-cohomology.p002
The classical notation for
$r_0\otimes\cdots\otimes r_{n+1}\in\mathsf{B}_nR$ is
$(r_0|\cdots|r_{n+1})$; this is the origin of the name ``bar.'' The
homomorphism $b$ acts by deleting one separator in every possible way and
then summing the results with alternating signs. From this it is easy to see
that $b^2=0$. Indeed, $b^2(r_0|\cdots|r_{n+1})$ is a linear combination of
elements of the form
% segment-id: o014.li2.chapter3.exercises-hochschild-homology-and-cohomology.q003
\begin{gather*}
	(\cdots |r_{h-1} r_h| \cdots |r_{k-1} r_k| \cdots )
	\quad \text{or} \quad
	( \cdots | r_{h-1} r_h r_{h+1} | \cdots ).
\end{gather*}
There are exactly two ways to obtain any such term from
$(\cdots|r_{h-1}|\cdots|r_k|\cdots)$ by deleting separators, and their signs
cancel. Thus $\mathsf{B}R$ is indeed a chain complex.

% segment-id: o014.li2.chapter3.exercises-hochschild-homology-and-cohomology.p003
Now augment the bar complex as in the following diagram and denote the result
by $\mathsf{B}'R$. The morphism $b_0$ in the diagram is also called the
augmentation homomorphism.
% segment-id: o014.li2.chapter3.exercises-hochschild-homology-and-cohomology.g001
\[\begin{tikzcd}[row sep=small]
	\cdots \arrow[r, "b_2"] & \mathsf{B}_1 R \arrow[r, "b_1"]
	& \mathsf{B}_0 R \arrow[r, "b_0"] & (\mathsf{B}_{-1} R := R) \arrow[r] & 0 \\
	& & (r_0 | r_1) \arrow[mapsto, r]
	\arrow[phantom, u, "\in" description, sloped]
	& r_0 r_1 \arrow[phantom, u, "\in" description, sloped] &
\end{tikzcd}\]\index[sym1]{B'R@$\mathsf{B}'R$}

% segment-id: o014.li2.chapter3.exercises-hochschild-homology-and-cohomology.le001
\begin{lemma}\label{prop:HH-bar-exactness}
	The augmented chain complex $\mathsf{B}'R$ is exact. More precisely, if
	$\mathsf{B}'R$ is regarded as a chain complex of $\Bbbk$-modules, then
	$\identity_{\mathsf{B}'R}$ is null-homotopic.
\end{lemma}

% segment-id: o014.li2.chapter3.exercises-hochschild-homology-and-cohomology.pf001
\begin{proof}
	We construct a family of $\Bbbk$-module homomorphisms
	$h_n:\mathsf{B}'_nR\to\mathsf{B}'_{n+1}R$, for $n\geq-1$, such that
	$b_{n+1}h_n+h_{n-1}b_n=\identity_{\mathsf{B}'_nR}$ (with, of course, the
	conventions $h_{-2}=0$ and $b_{-1}=0$). This will prove that
	$\identity_{\mathsf{B}'R}$ is null-homotopic. Concretely, take
% segment-id: o014.li2.chapter3.exercises-hochschild-homology-and-cohomology.q004
	\[ h_n(r_0 | \cdots | r_{n+1})
		= (1_R | r_0 | \cdots | r_{n+1}), \]
	and extend linearly on $\mathsf{B}'_nR$. For every $n\geq0$ and
	$r_0,\ldots,r_{n+1}\in R$, we obtain
% segment-id: o014.li2.chapter3.exercises-hochschild-homology-and-cohomology.q005
	\begin{align*}
		b_{n+1} h_n \left(r_0 | \cdots | r_{n+1} \right)
		& = (r_0 | \cdots | r_{n+1})
		+ \sum_{k=0}^n (-1)^{k+1}
		(1_R | \cdots | r_k r_{k+1} | \cdots), \\
		h_{n-1} b_n \left(r_0 | \cdots | r_{n+1} \right)
		& = \sum_{k=0}^n (-1)^k
		(1_R | \cdots | r_k r_{k+1} | \cdots).
	\end{align*}
	It follows immediately that
% segment-id: o014.li2.chapter3.exercises-hochschild-homology-and-cohomology.q006
	\[ (b_{n+1} h_n + h_{n-1} b_n)(r_0 | \cdots | r_{n+1})
		= (r_0 | \cdots | r_{n+1}). \]
	Since $b_0h_{-1}(r)=b_0(1_R|r)=r$, the equality also holds directly for
	$n=-1$.
\end{proof}

% segment-id: o014.li2.chapter3.exercises-hochschild-homology-and-cohomology.p004
At first sight, the definitions of $\mathsf{B}R$ and $\mathsf{B}'R$ may seem
like unexpected inventions. Example~\sourcecrossref{eg:HH-comonad}{in the
discussion of Chapter~8} will later view these chain complexes from the
perspective of comonads and give a natural explanation.

% segment-id: o014.li2.chapter3.exercises-hochschild-homology-and-cohomology.cn001
\begin{convention}\label{con:Re-R-R}
	\index[sym1]{R-e@$R^e$}
	Define $R^e:=R\otimes R^{\opp}$. For every $(R,R)$-bimodule $M$,
	including the special case $M=R$, henceforth:
% segment-id: o014.li2.chapter3.exercises-hochschild-homology-and-cohomology.l001
	\begin{compactitem}
% segment-id: o014.li2.chapter3.exercises-hochschild-homology-and-cohomology.i001
		\item regard $M$ as a left $R^e$-module through
		$(r\otimes r')m=rmr'$;
% segment-id: o014.li2.chapter3.exercises-hochschild-homology-and-cohomology.i002
		\item regard $M$ as a right $R^e$-module through
		$m(r\otimes r'):=r'mr$.
	\end{compactitem}
\end{convention}

% segment-id: o014.li2.chapter3.exercises-hochschild-homology-and-cohomology.p005
With this convention, for every $M$ we can define the chain complex of
$\Bbbk$-modules $M\dotimes{R^e}\mathsf{B}R$ and the complex
$\Hom_{R^e}\left(\mathsf{B}R,M\right)$.

% segment-id: o014.li2.chapter3.exercises-hochschild-homology-and-cohomology.de002
\begin{definition}[G.\ Hochschild]\label{def:HH}
	\index{Hochschild homology and cohomology}
	\index[sym1]{HH@$\HHm_n$, $\HHm^n$}
	For every $n$, define the $\Bbbk$-linear functors
	$\HHm_n,\HHm^n:(R,R)\dcate{Mod}\rightrightarrows\Bbbk\dcate{Mod}$ as
	follows:
% segment-id: o014.li2.chapter3.exercises-hochschild-homology-and-cohomology.q007
	\[\begin{array}{cl}
		\text{Hochschild homology}
		& \HHm_n(M) := \Hm_n\left(M\dotimes{R^e}\mathsf{B}R\right), \\
		\text{Hochschild cohomology}
		& \HHm^n(M) := \Hm^n\left(\Hom_{R^e}(\mathsf{B}R,M)\right).
	\end{array}\]
	In the special case $M=R$, these objects are called respectively the
	\emph{Hochschild homology} $\HHm_n(R)$ and \emph{Hochschild cohomology}
	$\HHm^n(R)$ of $R$.
\end{definition}

% segment-id: o014.li2.chapter3.exercises-hochschild-homology-and-cohomology.p006
To simplify the descriptions of $\HHm_n(M)$ and $\HHm^n(M)$, introduce the
following two Hochschild complexes:
% segment-id: o014.li2.chapter3.exercises-hochschild-homology-and-cohomology.q008
\begin{equation}\label{eqn:HH-cplx}\begin{aligned}
	\index[sym1]{CRM@$C_\bullet(R, M)$, $C^\bullet(R, M)$}
	C_\bullet(R,M) & := \left[
	\cdots \to M \otimes R^{\otimes n} \xrightarrow{d_n} \cdots
	\xrightarrow{d_2} M \otimes R \xrightarrow{d_1} M \right], \\
	C^\bullet(R,M) & := \left[
	M \xrightarrow{d^0} \Hom_{\Bbbk}(R, M) \xrightarrow{d^1} \cdots
	\to \Hom_{\Bbbk}(R^{\otimes n}, M) \xrightarrow{d^n} \cdots \right].
\end{aligned}\end{equation}
\index{Hochschild complex}
Their degrees are respectively $\ldots,2,1,0$ and $0,1,2,\ldots$; all other
terms are set equal to $0$. The homomorphisms $d_n$ and $d^n$ are defined as
follows.
% segment-id: o014.li2.chapter3.exercises-hochschild-homology-and-cohomology.l002
\begin{enumerate}
% segment-id: o014.li2.chapter3.exercises-hochschild-homology-and-cohomology.i003
	\item In bar notation, continue to write an element
	$m\otimes r_1\otimes\cdots\otimes r_n$ of $M\otimes R^{\otimes n}$ as
	$(m|r_1|\cdots|r_n)$. Define
% segment-id: o014.li2.chapter3.exercises-hochschild-homology-and-cohomology.q009
	\begin{multline}\label{eqn:HH-cplx-d}
		d_n(m|r_1|\cdots|r_n)=\\
		\underbracket{(mr_1|r_2|\cdots|r_n)
		+\sum_{k=1}^{n-1}(-1)^k
		(m|\cdots|r_kr_{k+1}|\cdots)}_{=:d'_n(m|r_1|\cdots|r_n)}
		+(-1)^n(r_nm|r_1|\cdots|r_{n-1}).
	\end{multline}
% segment-id: o014.li2.chapter3.exercises-hochschild-homology-and-cohomology.i004
	\item Identify an element of $\Hom_{\Bbbk}(R^{\otimes n},M)$ with a
	$\Bbbk$-multilinear map $R^n\to M$ (see \cite[\S 7.5]{Li1}). Define
% segment-id: o014.li2.chapter3.exercises-hochschild-homology-and-cohomology.q010
	\begin{multline*}
		(d^nf)(r_1,\ldots,r_{n+1})=\\
		r_1f(r_2,\ldots,r_{n+1})
		+\sum_{k=1}^n(-1)^kf(\ldots,r_kr_{k+1},\ldots)
		+(-1)^{n+1}f(r_1,\ldots,r_n)r_{n+1}.
	\end{multline*}
\end{enumerate}

% segment-id: o014.li2.chapter3.exercises-hochschild-homology-and-cohomology.p007
For every $n\in\Z_{\geq0}$, there are isomorphisms
% segment-id: o014.li2.chapter3.exercises-hochschild-homology-and-cohomology.g002
\[\begin{tikzcd}[row sep=tiny]
	M \dotimes{R^e} \mathsf{B}_n R \arrow[leftrightarrow, r, "\sim"]
	& M \otimes R^{\otimes n} \\
	m \otimes (r_0 \otimes \cdots \otimes r_{n+1}) \arrow[mapsto, r]
	& (r_{n+1} m r_0 | r_1 | \cdots | r_n) \\
	m \otimes (1_R \otimes r_1 \otimes \cdots \otimes r_n \otimes 1_R )
	& (m | r_1 | \cdots | r_n) \arrow[mapsto, l]
\end{tikzcd}\]
\footnote{Translator's note (O014-C044): the source appends a final
$|1_R$ to the image in the middle row, even though the codomain is
$M\otimes R^{\otimes n}$; that surplus component is removed here.}
and
% segment-id: o014.li2.chapter3.exercises-hochschild-homology-and-cohomology.g003
\[\begin{tikzcd}[row sep=tiny, column sep=small]
	\Hom_{R^e}(\mathsf{B}_nR, M) \arrow[leftrightarrow, r, "\sim"]
	& \Hom_{\Bbbk}(R^{\otimes n}, M) \\
	\varphi \arrow[mapsto, r]
	& {\left[(r_1,\ldots,r_n)\mapsto
		\varphi(1_R,r_1,\ldots,r_n,1_R)\right]} \\
	{\left[r_0\otimes\cdots\otimes r_{n+1}\mapsto
		r_0f(r_1,\ldots,r_n)r_{n+1}\right]}
	& f \arrow[mapsto, l] .
\end{tikzcd}\]
A short calculation shows that, through these isomorphisms,
$\identity_M\otimes b_n$ corresponds to $d_n$, while $b_{n+1}^*$
corresponds to $d^n$.\footnote{Translator's note (O014-C045): the source
writes $b_n^*$, which starts at $\Hom_{R^e}(\mathsf{B}_{n-1}R,M)$; the
degree-$n$ differential comes from
$b_{n+1}:\mathsf{B}_{n+1}R\to\mathsf{B}_nR$.}
Thus the isomorphisms above are isomorphisms of complexes and give the
following result. To display the parameter $R$ explicitly once again, write
$\HHm_n(R,M):=\HHm_n(M)$ and $\HHm^n(R,M):=\HHm^n(M)$.
\footnote{Translator's note (O014-C046): in the formula below, the source
switches to two-argument notation without defining it; these two
abbreviations make that transition explicit.} With this notation,
% segment-id: o014.li2.chapter3.exercises-hochschild-homology-and-cohomology.q011
\[ \HHm^n(R, M) \simeq \Hm^n(C^\bullet(R, M)), \quad
	\HHm_n(R, M) \simeq \Hm_n(C_\bullet(R, M)). \]
Readers new to the subject are encouraged to check the calculation in
detail.

% segment-id: o014.li2.chapter3.exercises-hochschild-homology-and-cohomology.re001
\begin{remark}
	Suppose that $R$ is commutative. Every $R$-module $M$ becomes an
	$(R,R)$-bimodule by $rmr':=rr'm$. In this case, the complex
	$C_\bullet(R,M)$ (or $C^\bullet(R,M)$)\footnote{Translator's note
	(O014-C047): the source reverses the two arguments as $C_\bullet(M,R)$
	and $C^\bullet(M,R)$ in this sentence; the notation here is aligned with
	the preceding definitions.}
	is a complex of $R$-modules if we set
% segment-id: o014.li2.chapter3.exercises-hochschild-homology-and-cohomology.q012
	\[ r\cdot(m|r_1|\cdots|r_n)=(rm|r_1|\cdots|r_n)
		\quad\text{or}\quad
		(r\cdot f)(r_1,\ldots,r_n)=rf(r_1,\ldots,r_n). \]
	Consequently, $\HHm_n(M)$ and $\HHm^n(M)$ both acquire $R$-module
	structures. This can of course also be proved at the level of the bar
	complex.
\end{remark}

% segment-id: o014.li2.chapter3.exercises-hochschild-homology-and-cohomology.ex001
\begin{example}
	Take $M=R=\Bbbk$. It follows immediately from the construction above that
	$d_1,d_2,\ldots$ are successively $0,\identity,0,\identity,\ldots$.
	Thus $\HHm_0(\Bbbk)=\Bbbk$ and $\HHm_{\geq1}(\Bbbk)=0$. A similar
	argument gives $\HHm^0(\Bbbk)=\Bbbk$ and
	$\HHm^{\geq1}(\Bbbk)=0$.
\end{example}

% segment-id: o014.li2.chapter3.exercises-hochschild-homology-and-cohomology.p008
In higher degrees $n$, computing $\HHm^n(M)$ and $\HHm_n(M)$ from the
original definition is generally difficult.
Example~\sourcecrossref{eg:HH-Tor}{in \S~3.14} and
\sourcecrossref{eg:HH-Ext}{in \S~3.14} will later show respectively that
% segment-id: o014.li2.chapter3.exercises-hochschild-homology-and-cohomology.q013
\[ \HHm_n(M) \simeq \Tor^{R^e}_n(M, R), \quad
	\HHm^n(M) \simeq \Ext_{R^e}^n(R, M), \]
provided that $R$, as a $\Bbbk$-module, is respectively flat and projective.
At that point, some special cases of $\HHm_n$ and $\HHm^n$ can be computed
with simpler complexes. For general $R$, in the exercises of
\sourcecrossref{sec:simplicial}{Chapter~8} we will explain $\HHm_n$ and
$\HHm^n$ as relative $\Tor$ and relative $\Ext$.

% segment-id: o014.li2.chapter3.exercises-hochschild-homology-and-cohomology.ex002
\begin{example}[Degree zero: center and cocenter]
	\index{center}
	\index{cocenter}
	First consider $\HHm^0(M)$. By definition,
	$d^0:M=C^0(R,M)\to C^1(R,M)=\Hom_{\Bbbk}(R,M)$ maps $m$ to
	$[r\mapsto rm-mr]$. Therefore
% segment-id: o014.li2.chapter3.exercises-hochschild-homology-and-cohomology.q014
	\[ \HHm^0(M)=\left\{m\in M:\forall r\in R,\;rm=mr\right\}. \]
	The right-hand side may naturally be called the \emph{center} of $M$;
	when $M=R$, this is precisely the center in the sense of ring theory.

	Next consider $\HHm_0(M)$. Write $[M,R]$ for the $\Bbbk$-submodule of
	$M$ generated by all elements of the form $mr-rm$, with $m\in M$ and
	$r\in R$. Since $d_1(m|r)=mr-rm$, we obtain
% segment-id: o014.li2.chapter3.exercises-hochschild-homology-and-cohomology.q015
	\[ \Image\left[d_1:M\otimes R\to M\right]=[M,R],
		\quad \HHm_0(M)=M/[M,R]. \]
	In particular, the case $M=R$ gives the submodule $[R,R]$ of $R$
	generated by all $[r',r]:=r'r-rr'$. The corresponding $\Bbbk$-module
	$R/[R,R]$ is called the \emph{cocenter} of $R$. Every $\Bbbk$-module
	homomorphism $\varphi:R\to N$ satisfying $\varphi(rr')=\varphi(r'r)$,
	that is, having a ``trace-like'' property, factors uniquely through
	$R/[R,R]=\HHm_0(R)$.
\end{example}

% segment-id: o014.li2.chapter3.exercises-hochschild-homology-and-cohomology.ex003
\begin{example}[Degree one: derivations]\label{eg:HH1}
	\index[sym1]{Der@$\mathrm{Der}_{\Bbbk}(R, M)$}
	\index[sym1]{InnDer@$\mathrm{Inn}_{\Bbbk}(R, M)$}
	Now consider $\HHm^1(M)$. Write $[r,m]:=rm-mr$. Then
% segment-id: o014.li2.chapter3.exercises-hochschild-homology-and-cohomology.q016
	\begin{align*}
		\Ker(d^1)
		&=\left\{D\in\Hom_{\Bbbk}(R,M):
		\forall r_1,r_2\in R,\;
		r_1D(r_2)-D(r_1r_2)+D(r_1)r_2=0\right\},\\
		\Image(d^0)
		&=\left\{[\mathord{\cdot},m]\in\Hom_{\Bbbk}(R,M):m\in M\right\}.
	\end{align*}
	The condition defining $\Ker(d^1)$ may be rewritten as the Leibniz rule
	$D(r_1r_2)=r_1D(r_2)+D(r_1)r_2$. A $\Bbbk$-module homomorphism $D$ with
	this property is regarded as a derivation on $R$ with values in $M$. The
	$\Bbbk$-module formed by these homomorphisms is denoted by
	$\Der_{\Bbbk}(R,M)$, while the elements of the form
	$[\mathord{\cdot},m]$ form the submodule
	$\mathrm{Inn}_{\Bbbk}(R,M)$. Thus the quotient module
% segment-id: o014.li2.chapter3.exercises-hochschild-homology-and-cohomology.q017
	\[ \HHm^1(M)=\Der_{\Bbbk}(R,M)/\mathrm{Inn}_{\Bbbk}(R,M) \]
	classifies all derivations on $R$ with values in $M$, modulo
	$\mathrm{Inn}_{\Bbbk}(R,M)$.

	Now suppose that $R$ is commutative, in order to interpret $\HHm_1(M)$.
	We need a little preparation. Define the free $R$-module
	$\bigoplus_{r\in R}R\widetilde{\dd r}$ with basis the symbols
	$\widetilde{\dd r}$, and define the submodule $N$ generated by the
	following elements:
% segment-id: o014.li2.chapter3.exercises-hochschild-homology-and-cohomology.q018
	\[ \widetilde{\dd(r+r')}-\widetilde{\dd r}-\widetilde{\dd r'},
		\quad \widetilde{\dd tr}-t\widetilde{\dd r},
		\quad \widetilde{\dd rr'}-r\widetilde{\dd r'}-r'\widetilde{\dd r}, \]
	where $r,r'\in R$ and $t\in\Bbbk$. Thus define the $R$-module
% segment-id: o014.li2.chapter3.exercises-hochschild-homology-and-cohomology.q019
	\begin{align*}
		\Omega_{R|\Bbbk}
		&:=\bigoplus_{r\in R}R\widetilde{\dd r}\bigg/N\\
		&=\sum_{r\in R}R\dd r,\qquad
		\dd r:=\text{the image of }\widetilde{\dd r}.
	\end{align*}
	\index{Kähler differentials}
	\index[sym1]{OmegaRk@$\Omega_{R\vert\Bbbk}$}
	This is called the \emph{module of Kähler differentials} of the
	$\Bbbk$-algebra $R$ and is characterized by the following universal
	property:
% segment-id: o014.li2.chapter3.exercises-hochschild-homology-and-cohomology.g004
	\[\begin{tikzcd}[row sep=tiny]
		\Hom_R(\Omega_{R|\Bbbk}, M) \arrow[leftrightarrow, r, "\sim"]
		& \Der_{\Bbbk}(R, M) \\
		\varphi \arrow[mapsto, r] & {\left[r\mapsto\varphi(\dd r)\right]} \\
		{\left[\dd r\mapsto D(r)\right]} & D \arrow[mapsto, l].
	\end{tikzcd}
	\qquad M:\text{an $R$-module}. \]
	The reader is asked to verify this directly. This isomorphism shows that
	$R\to\Omega_{R|\Bbbk}$ given by $r\mapsto\dd r$ (corresponding to
	$\varphi=\identity$) is a ``universal derivation.'' Topics related to
	this properly belong to the theory of commutative rings; we merely record
	the result here.

	Return to $\HHm_1(M)$, still assuming that $R$ is commutative. Let $M$ be
	an $R$-module made into an $(R,R)$-bimodule by $rmr':=(rr')m$. It is
	immediate that $d_1:M\otimes R\to M$ is $0$, while the image of
	$d_2:M\otimes R^{\otimes2}\to M\otimes R$ is generated by elements of
	the form $(rm|r')-(m|rr')+(r'm|r)$. We therefore obtain homomorphisms of
	$\Bbbk$-modules in both directions
% segment-id: o014.li2.chapter3.exercises-hochschild-homology-and-cohomology.g005
	\[\begin{tikzcd}[row sep=tiny]
		(M \otimes R) / \Image(d_2) \arrow[shift right, r]
		& M \dotimes{R} \Omega_{R|\Bbbk} \arrow[shift right, l] \\
		(m|r) + \Image(d_2) \arrow[mapsto, r] & m \otimes \dd r \\
		(r'm|r) + \Image(d_2)
		& m \otimes r' \dd r \arrow[mapsto, l] .
	\end{tikzcd}\]
	From the descriptions of $\Image(d_2)$ and $\Omega_{R|\Bbbk}$, a routine
	check shows that these homomorphisms are well defined and mutually
	inverse. In fact, both are isomorphisms of $R$-modules. Thus, when $R$ is
	commutative, we obtain
	$\HHm_1(M)\simeq M\dotimes{R}\Omega_{R|\Bbbk}$. In particular,
	$\HHm_1(R)\simeq\Omega_{R|\Bbbk}$.
\end{example}

% segment-id: o014.li2.chapter3.exercises-hochschild-homology-and-cohomology.p009
Hochschild homology and cohomology contain a wealth of information; their
connection with derivations and differential forms is no accident. The
exercises will give further interpretations of Hochschild cohomology.

% segment-id: o014.li2.chapter3.exercises-hochschild-homology-and-cohomology.p010
Return to the Hochschild chain complex. Keeping the actions of $R^e$ on
$\mathsf{B}_nR$ and on $M$ in mind, the intuitive idea is to arrange an
element $(m|r_1|\cdots|r_n)$ of $C_n(R,M)$ in a circle:
% segment-id: o014.li2.chapter3.exercises-hochschild-homology-and-cohomology.f001
\begin{center}\begin{tikzpicture}[scale=1.2]
		\node[rotate=30] at (120:1) {$r_n$};
		\node[rotate = 0]  at (90:1) {$m$};
		\node[rotate=-15] at (60:1) {$r_1$};
		\node[rotate=-45]  at (30:1) {$r_2$};
		\draw (135:0.8) -- (135:1.2);
		\draw (105:0.8) -- (105:1.2);
		\draw (75:0.8) -- (75:1.2);
		\draw(45:0.8) -- (45:1.2);
		\draw[loosely dotted, thick]
		(10:1) arc [start angle=10, end angle=-210, radius=1];
\end{tikzpicture}\end{center}
Thus $d_n(m|r_1|\cdots|r_n)$ is obtained by deleting a separator in $n+1$
ways and summing the results with alternating signs. In the special case
$M=R$, the picture has an evident rotational symmetry. Algebraically, for
every $n\in\Z_{\geq0}$ define the $\Bbbk$-module homomorphism
% segment-id: o014.li2.chapter3.exercises-hochschild-homology-and-cohomology.g006
\[\begin{tikzcd}[row sep=tiny]
	t_n: R^{\otimes (n+1)} \arrow[r] & R^{\otimes (n+1)} \\
	(r_0 | \cdots | r_n) \arrow[mapsto, r]
	& (-1)^n \cdot (r_n | r_0 | \cdots | r_{n-1}) ;
\end{tikzcd}\]
then $t_n^{n+1}=\identity$. This homomorphism realizes the symmetry.

% segment-id: o014.li2.chapter3.exercises-hochschild-homology-and-cohomology.p011
Rotational symmetry leads to the theory of \emph{cyclic homology}.
Historically, there have been at least two motivations for studying cyclic
homology. The first is to seek an analogue of de Rham theory in the
noncommutative setting. The second is to study and apply K-theory, including
various generalizations of the index theorem. This section uses Hochschild
homology and cohomology and cyclic homology only as a light vehicle for
familiarizing the reader with operations on complexes; the treatment is
merely introductory. Readers wishing to go further may consult monographs
such as \cite{Lo98, Wi19}.

% segment-id: o014.li2.chapter3.exercises-hochschild-homology-and-cohomology.p012
All complexes (or double complexes) mentioned below are understood to be
chain complexes (or chain double complexes).

% segment-id: o014.li2.chapter3.exercises-hochschild-homology-and-cohomology.p013
For every $n\geq0$, define
$N_n:=\identity+t_n+\cdots+(t_n)^n
\in\End_{\Bbbk}(R^{\otimes(n+1)})$. Define the \emph{cyclic double complex}
$\mathrm{CC}(R)=(\mathrm{CC}(R)_{p,q})_{(p,q)\in\Z^2}$ to be the following
double complex.
\index{cyclic double complex}
\index[sym1]{CCR@$\mathrm{CC}(R)$}
% segment-id: o014.li2.chapter3.exercises-hochschild-homology-and-cohomology.g007
\[\begin{tikzcd}[
	/tikz/execute at end picture={
		\node[rectangle, draw, fit=(NW) (SE),
		inner xsep=1.5em, inner ysep=1.5em] {};
	}]
	\vdots & |[alias=NW]| & \vdots \arrow[d, "{d_{q+1}}"]
	& \vdots \arrow[d, "{d'_{q+1}}"] & \\
	q & \cdots & R^{\otimes (q+1)} \arrow[d, "{d_q}"] \arrow[l, "{N_q}"]
	& R^{\otimes (q+1)} \arrow[d, "{d'_q}"]
	\arrow[l, "{\identity - t_q}"] & \cdots \arrow[l] \\
	q-1 & \cdots & R^{\otimes q} \arrow[d] \arrow[l, "{N_{q-1}}"]
	& R^{\otimes q} \arrow[d] \arrow[l, "{\identity - t_{q-1}}"]
	& \cdots \arrow[l] \\
	\vdots & & \vdots & \vdots & |[alias=SE]| \\
	& \cdots & p: \text{even} & p+1: \text{odd} & \cdots
\end{tikzcd}\]
All its terms are $\Bbbk$-modules, and all terms with $q<0$ are defined to
be $0$. The morphisms $d_q$ and $d'_q$ are defined in
\eqref{eqn:HH-cplx-d}. Notice that:
% segment-id: o014.li2.chapter3.exercises-hochschild-homology-and-cohomology.l003
\begin{compactitem}
% segment-id: o014.li2.chapter3.exercises-hochschild-homology-and-cohomology.i005
	\item the even columns are the Hochschild chain complex $C_\bullet(R,R)$
	of \eqref{eqn:HH-cplx};
% segment-id: o014.li2.chapter3.exercises-hochschild-homology-and-cohomology.i006
	\item it is easy to prove that the odd columns are the augmented bar
	complex $\mathsf{B}'R$, with the indices shifted so that it begins in
	degree $0$. Hence Lemma~\ref{prop:HH-bar-exactness} implies that all odd
	columns are exact; in fact, they are $0$ in
	$\cate{K}(\Bbbk\dcate{Mod})$.
\end{compactitem}

% segment-id: o014.li2.chapter3.exercises-hochschild-homology-and-cohomology.p014
To show that $\mathrm{CC}(R)$ is indeed a double complex, we need the
following observation. Its proof is elementary and interesting, with no
essential difficulty, so it is left as an exercise in this chapter.

% segment-id: o014.li2.chapter3.exercises-hochschild-homology-and-cohomology.le002
\begin{lemma}\label{prop:CC-cplx}
	Let $R$ be a $\Bbbk$-algebra. For every $q\in\Z_{\geq1}$, the following
	equalities hold in
	$\Hom_{\Bbbk}\left(R^{\otimes(q+1)},R^{\otimes q}\right)$:
% segment-id: o014.li2.chapter3.exercises-hochschild-homology-and-cohomology.q020
	\[ d_q(\identity-t_q)=(\identity-t_{q-1})d'_q,
		\quad d'_qN_q=N_{q-1}d_q. \]
\end{lemma}

% segment-id: o014.li2.chapter3.exercises-hochschild-homology-and-cohomology.p015
For every double complex $C=(C_{i,j})_{(i,j)\in\Z^2}$ and $m\in\Z$,
define its horizontal shift\footnote{This book's convention uses the
subscript $\mathrm{I}$ for horizontal operations and $\mathrm{II}$ for
vertical operations.} $C_{\mathrm{I}}[m]$ to be the double complex
$(C_{m+i,j})_{i,j}$. For every $p\in\Z$, define the horizontal brutal
truncation functor $\sigma_{\mathrm{I},\leq p}C$ (or
$\sigma_{\mathrm{I},\geq p}C$) by replacing with $0$ every $C_{i,j}$ for
which $i>p$ (or $i<p$) and leaving all other terms unchanged. This gives a
short exact sequence of double complexes (note the order):
% segment-id: o014.li2.chapter3.exercises-hochschild-homology-and-cohomology.q021
\begin{equation}\label{eqn:homology-I-brute-truncation}
	0 \to \sigma_{\mathrm{I}, \leq p} C \to C
	\to \sigma_{\mathrm{I}, \geq p+1} C \to 0.
\end{equation}
Moreover, for every $a\leq b$, define
$\sigma_{\mathrm{I},[a,b]}:=
\sigma_{\mathrm{I},\leq b}\sigma_{\mathrm{I},\geq a}
=\sigma_{\mathrm{I},\geq a}\sigma_{\mathrm{I},\leq b}$.


% segment-id: o014.li2.chapter3.exercises-hochschild-homology-and-cohomology.p016
Apply these functors to $\mathrm{CC}(R)$ and then take the total complex
$\tot_{\Pi}$ (Definition~\ref{def:total-cplx}); this leads to the following
definition.

% segment-id: o014.li2.chapter3.exercises-hochschild-homology-and-cohomology.de003
\begin{definition}[B.\ Feigin, B.\ Tsygan; A.\ Connes]\label{def:HC}
	\index{cyclic homology}
	\index{periodic cyclic homology}
	\index[sym1]{HPHC@$\mathrm{HP}_n$, $\mathrm{HC}_n$, $\mathrm{HP}^n$, $\mathrm{HC}^n$}
	For a $\Bbbk$-algebra $R$ and every $n\in\Z$, define
% segment-id: o014.li2.chapter3.exercises-hochschild-homology-and-cohomology.q022
	\begin{align*}
		\mathrm{HP}_n(R)
		&:=\Hm_n\left(\tot_{\Pi}(\mathrm{CC}(R))\right)
		&\text{(periodic cyclic homology)},\\
		\mathrm{HC}_n(R)
		&:=\Hm_n\left(
		\tot\left(\sigma_{\mathrm{I},\geq0}\mathrm{CC}(R)\right)\right)
		&\text{(cyclic homology)}.
%		\mathrm{HC}^-_n(R) & := \Hm_n\left( \tot_{\Pi}\left(\sigma_{\mathrm{I}, \leq 1} \mathrm{CC}(R)\right) \right) & \text{(negative cyclic homology)}
	\end{align*}
	Both are functors $\Bbbk\dcate{Alg}\to\Bbbk\dcate{Mod}$.
\end{definition}

% segment-id: o014.li2.chapter3.exercises-hochschild-homology-and-cohomology.p017
Since $\sigma_{\mathrm{I},\geq0}\mathrm{CC}(R)$ lies in the first quadrant,
its total complex involves only the finite direct sums
$\bigoplus_{\substack{p+q=n\\p,q\geq0}}\mathrm{CC}_{p,q}(R)$. The subscript
on $\tot$ is therefore unnecessary.

% segment-id: o014.li2.chapter3.exercises-hochschild-homology-and-cohomology.p018
The definition immediately gives $\mathrm{HC}_{<0}(R)=0$, while
$\mathrm{HC}_0(R)$ is the quotient of $R$ by
$\Image[d_1:R^{\otimes2}\to R]$ and $\Image[\identity-t_0]=0$, namely
$R/[R,R]$. The exercises will give computations in more special cases.

% segment-id: o014.li2.chapter3.exercises-hochschild-homology-and-cohomology.p019
The cyclic complex has the periodicity
$\mathrm{CC}(R)=\mathrm{CC}(R)_{\mathrm{I}}[-2]$, which yields an
isomorphism $\mathrm{HP}_n(R)\rightiso\mathrm{HP}_{n-2}(R)$. For cyclic
homology, the corresponding construction is the epimorphism ``shift two
columns to the left''
% segment-id: o014.li2.chapter3.exercises-hochschild-homology-and-cohomology.q023
\begin{align*}
	\sigma_{\mathrm{I},\geq0}\mathrm{CC}(R)
	&\to\left(\sigma_{\mathrm{I},\geq0}\mathrm{CC}(R)\right)_{\mathrm{I}}[-2]\\
	\mathrm{CC}(R)_{p,q}
	&\to\begin{cases}
		\mathrm{CC}(R)_{p-2,q}\;\text{(identity)},&p\geq2,\\
		0,&0\leq p<2.
	\end{cases}
\end{align*}
Its kernel is given by columns $0$ and $1$ of $\mathrm{CC}(R)$, namely the
sub-double-complex $\sigma_{\mathrm{I},[0,1]}\mathrm{CC}(R)$. This defines
Connes's periodicity operator
$S:\mathrm{HC}_n(R)\to\mathrm{HC}_{n-2}(R)$ on cyclic homology.
\index{Connes periodicity operator}

% segment-id: o014.li2.chapter3.exercises-hochschild-homology-and-cohomology.th001
\begin{theorem}[A.\ Connes]\label{prop:Connes-exact}
	Let $R$ be a $\Bbbk$-algebra. There is a canonical long exact sequence
% segment-id: o014.li2.chapter3.exercises-hochschild-homology-and-cohomology.q024
	\[ \cdots \xrightarrow{S} \mathrm{HC}_{n-1}(R)
		\xrightarrow{B} \HHm_n(R) \xrightarrow{I} \mathrm{HC}_n(R)
		\xrightarrow{S} \mathrm{HC}_{n-2}(R) \to \cdots, \]
	where $S$ is Connes's periodicity operator and $I$ is induced by the
	inclusion
	\[
		C_\bullet(R,R)\xrightarrow{\text{column $0$}}
		\sigma_{\mathrm{I},\geq0}\mathrm{CC}(R).
	\]
	The map $B$ is the corresponding connecting homomorphism.
\end{theorem}

% segment-id: o014.li2.chapter3.exercises-hochschild-homology-and-cohomology.pf002
\begin{proof}
	The first step is to apply \eqref{eqn:homology-I-brute-truncation} to
	obtain the short exact sequence of double complexes
% segment-id: o014.li2.chapter3.exercises-hochschild-homology-and-cohomology.g008
	\[\begin{tikzcd}
		0 \arrow[r] & \text{column $0$} \arrow[r]
		& \sigma_{\mathrm{I},[0,1]}\mathrm{CC}(R) \arrow[r]
		& \text{column $1$} \arrow[r] & 0 .
	\end{tikzcd}\]
	Taking total complexes involves only finite direct sums. Checking degree
	by degree shows that the result remains a short exact sequence
% segment-id: o014.li2.chapter3.exercises-hochschild-homology-and-cohomology.g009
	\[\begin{tikzcd}
		0 \arrow[r] & \text{column $0$} \arrow[r]
		& \tot\left(\sigma_{\mathrm{I},[0,1]}\mathrm{CC}(R)\right)
		\arrow[r] & \text{column $1$} \arrow[r] & 0.\\
		& C_\bullet(R,R) \arrow[u,"\sim" sloped]
		& & \mathsf{B}'R \arrow[u,"\sim" sloped] &
	\end{tikzcd}\]
	We already know that $\mathsf{B}'R$ is exact. Applying the long exact
	sequence of Proposition~\ref{prop:long-exact-sequence-ses} shows that
	$C_\bullet(R,R)\to
	\tot\left(\sigma_{\mathrm{I},[0,1]}\mathrm{CC}(R)\right)$ is a
	quasi-isomorphism of complexes.

	Next consider the short exact sequence mentioned when defining $S$:
% segment-id: o014.li2.chapter3.exercises-hochschild-homology-and-cohomology.g010
	\[\begin{tikzcd}
		0 \arrow[r] & \sigma_{\mathrm{I},[0,1]}\mathrm{CC}(R) \arrow[r]
		& \sigma_{\mathrm{I},\geq0}\mathrm{CC}(R) \arrow[r]
		& \left(\sigma_{\mathrm{I},\geq0}\mathrm{CC}(R)\right)_{\mathrm{I}}[-2]
		\arrow[r] & 0 .
	\end{tikzcd}\]
	After taking total complexes, this remains a short exact sequence. By the
	preceding step, the degree-$n$ homology groups of the three total complexes
	are identified respectively with $\HHm_n(R)$, $\mathrm{HC}_n(R)$, and
	$\mathrm{HC}_{n-2}(R)$.\footnote{Translator's note (O014-C048): the
	source says ``degree-$n$ cohomology'' in this sentence, but all three
	objects are chain complexes and the formulas use homological subscripts;
	``homology'' is used here.} This proves the theorem.
\end{proof}

% segment-id: o014.li2.chapter3.exercises-hochschild-homology-and-cohomology.re002
\begin{remark}
	\index{cyclic double complex}
	\index{periodic cyclic cohomology}
	\index{cyclic cohomology}
	\index[sym1]{HPHC}
	Define the cyclic double complex in the cochain sense,
	$\mathrm{CC}'(R)$, by
	$\mathrm{CC}'(R)^{p,q}:=
	\Hom_{\Bbbk}\left(\mathrm{CC}(R)_{p,q},\Bbbk\right)$. In analogy with
	Definition~\ref{def:HC}, define
% segment-id: o014.li2.chapter3.exercises-hochschild-homology-and-cohomology.q025
	\begin{align*}
		\mathrm{HP}^n(R)
		&:=\Hm^n\left(\tot_{\oplus}(\mathrm{CC}'(R))\right)
		&\text{(periodic cyclic cohomology)},\\
		\mathrm{HC}^n(R)
		&:=\Hm^n\left(
		\tot\left(\sigma_{\mathrm{I}}^{\geq0}\mathrm{CC}'(R)\right)\right)
		&\text{(cyclic cohomology)}
	\end{align*}
	and so forth, where $\sigma_{\mathrm{I}}^{\geq0}$ still denotes the
	horizontal brutal-truncation functor.

	Write $R^\vee:=\Hom_{\Bbbk}(R,\Bbbk)$ and make it an $(R,R)$-bimodule by
	\[ (r\varphi r')(x):=\varphi(r'xr). \]
	By currying, the $\Bbbk$-linear dual of the Hochschild column is identified
	with $C^\bullet(R,R^\vee)$; this identification requires no finiteness
	hypothesis. Consequently, the long exact sequence in
	Theorem~\ref{prop:Connes-exact} has a cohomological version
% segment-id: o014.li2.chapter3.exercises-hochschild-homology-and-cohomology.q026
	\[ \cdots \xrightarrow{S} \mathrm{HC}^{n+1}(R)
		\xrightarrow{I} \HHm^{n+1}(R,R^\vee) \xrightarrow{B}
		\mathrm{HC}^n(R) \xrightarrow{S} \mathrm{HC}^{n+2}(R) \to \cdots. \]
	\footnote{Translator's note (O014-C049): the source prints
	$\HHm^{n+1}(R)=\HHm^{n+1}(R,R)$ as the middle term. However, the
	$\Bbbk$-linear dual of $C_\bullet(R,R)$ has coefficients in $R^\vee$,
	not in the regular bimodule $R$, unless an additional identification
	$R^\vee\simeq R$ is supplied.}
	The corresponding Connes periodicity operator $S$ comes from the
	monomorphism ``shift two columns to the right,'' namely
% segment-id: o014.li2.chapter3.exercises-hochschild-homology-and-cohomology.q027
	\[ \sigma_{\mathrm{I}}^{\geq0}\mathrm{CC}'(R)
		\to\left(\sigma_{\mathrm{I}}^{\geq0}\mathrm{CC}'(R)\right)_{\mathrm{I}}[2]. \]
\end{remark}

% segment-id: o014.li2.chapter3.exercises-hochschild-homology-and-cohomology.p020
As our toolkit grows, we will return repeatedly to Hochschild homology and
cohomology and to cyclic homology and cohomology.
